\clearpage{\thispagestyle{empty}}
\section{Gröbner Bases}
With affine varieties under our belt, there is an obvious follow up to this knowledge. Especially for those interested in the real application of this mathematics. Namely, how do we solve these varieties? It is great to describe them, and on a theoretial level it is powerful but to apply them is a challenge in itself. Luckily, the problem of asking for the affine varity points boils down to solving a system of polynomials equations. This is easier said than done, and an issue arises when going to greater and greater dimensions. Most real world problems will involve many dimensions, and being able to solve these problems quickly and with confidence is of importance. To do so, a computer algebra system is most helpful.

Every computer language has important data structures that determine how the language is shaped. For Macualay 2, ideals is one of these types. They are important to algebra, and native support for manipulating ideals exist. However, to efficiently handle ideals, there is a problem that needs to be solved with great confidence. Namely: how do we go between abstract entities such as ideals and something that a computer can interact with?

\subsection{Dicksons Lemma}
The ideal description problem formulates the bridge between the abstract and machine world. It does this by concretizing the abstract concept of an ideal to a tangible set of generating polynomials that can be interacted with and evaluated. In order to bridge this gap, finite polynomials have to be found. Starting off with monomial ideals is the best way to do so. 

Dicksons Lemma 
\begin{definition}
Another monomial lies in $I = <x^a |  a \in A>$ if and only if $x^b$ is divisible by $x^a$ for some a in A.

\end{definition}
Proof missing.

Is needed to use in the Hilbert basis theorem. Monomial ideals have finite basis. Proof missing.


\subsection{Hilbert Bases}
Theorem 4 (Hilbert Basis Theorem). Every ideal has a finite generating set. Proof missing.

\subsection{Gröbner Bases}
Gröbner basis: 
$<LT(g_n)> = <LT(I)>$
Every ideal has a gröber basis and every basis of an ideal I is a basis of I. Proof and fluff missing.

\subsection{Properties of Gröbner Bases}
Properties of basis
Unique remainder when doing division which solves the issue of different solutions for different orders of the polynomial.